\documentclass[../main.tex]{subfiles}

\begin{document}
\section{\textbf{Cluster Expansion}}
\makebox[\textwidth][s]{Author: Xianchao \hfill Date: \today}

The cluster expansion is an approach used to construct a basis set. More precisely, it is discrete Fourier transformation on the configuration space.
On the basis, we can expand any lattice dependent only functions. It's a powerful tool for the investigation of atomic ordering.

Assuming we have a configuration space, $\Omega$. The single site in the configuration space is denoted as 
$\sigma_i$, and the configuration space is $(\sigma_1, \sigma_2, \dots, \sigma_N)$. For each site, it has different occupation
element, that is $A = \{\alpha_1, \alpha_2, \dots, \alpha_n\}$.

To define the orthonormal basis here, we have to give out the definition of the inner product first. The inner product
on the single site $\sigma_i$ is defined as:
\begin{equation}
    \braket{f}{g} =\sum_{a=1}^{n} \nu_a f(a)g(a)
\end{equation}
Here $\bm{\nu}$ is a reference distribution, which gives the probability of each specific element occupy this site.
The orthonormal basis is a set of basis that satisfy:
\begin{equation}
    \braket{\phi_m}{\phi_n} = \delta_{mn}
\end{equation}
It can be proved that, the number of the basis function is equal to the total number of the species.
The first basis function will always be constant:
\begin{equation}
    \phi_0 = 1
\end{equation}
The many body basis is defined as:
\begin{equation}
    \Psi_{S,\bm{r}}(\bm{\sigma}) = \prod_{i \in S}\phi_{r_i}(\sigma_i) 
\end{equation}
$S$ is the set of the site index in this many body term. $r_i \in \{1, 2, \dots, n-1\}$.

\end{document}
